\section{Avoidance Source Diagnosis}\label{sec:avoidance_source}

Fine-tuning on collision-free demonstrations produces genuine collision avoidance: LoRA-adapted Pi0.5 raises CAR by ${\sim}$26 percentage points with zero action-collapsed episodes. Does SDF geometric supervision contribute additional avoidance beyond what demonstrations provide? We compare LoRA-only against LoRA+SDF across a 20$\times$ range of loss weights ($\lambda \in \{0.005, 0.01, 0.1\}$) on Pi0.5, all evaluated at the 30K-step checkpoint.

\begin{table}[t]
\centering
\caption{Avoidance source attribution on SafeLIBERO. SDF shows no statistically significant CAR improvement beyond LoRA-only at the tested scale ($p > 0.05$ for $\lambda \in \{0.01, 0.1\}$). All fine-tuned models exhibit 0\% action-collapse rate and displacement $>$\,0.37\,m.}
\label{tab:avoidance_source}
\vspace{2pt}
\footnotesize
\begin{tabular*}{\columnwidth}{l@{\extracolsep{\fill}}cccc}
\toprule
Method & Seeds & TSR (\%) & CAR (\%) & Disp (m) \\
\midrule
Pi0.5 baseline & 1 & 71.5 & 28.5 & 0.367 \\
\midrule
LoRA-only & 3 & 49.7$\pm$3.2 & 54.83$\pm$2.25 & 0.391 \\
+SDF $\lambda{=}$0.005 & 3 & 54.0$\pm$3.0 & 54.5$\pm$2.29 & 0.387 \\
+SDF $\lambda{=}$0.01 & 3 & 53.8$\pm$1.5 & 56.33$\pm$1.61 & 0.389 \\
+SDF $\lambda{=}$0.1 & 3 & 50.2$\pm$3.9 & 57.00$\pm$2.18 & 0.394 \\
\bottomrule
\end{tabular*}
\vspace{1pt}
\end{table}

Table~\ref{tab:avoidance_source} shows that collision-free demonstrations alone raise CAR by ${\sim}$26 percentage points (28.5\% $\to$ 54.83\%) with zero action-collapsed episodes and displacement above 0.37\,m, confirming genuine avoidance. Adding SDF supervision does not yield statistically significant improvement at the tested scale: LoRA-only CAR (54.83\%) is statistically indistinguishable from all SDF conditions including $\lambda = 0.1$ (57.00\%, $p = 0.155$; Appendix~\ref{sec:appendix_null}). A Phase~2 ablation ($\lambda = 0$ after step 5K) yields CAR\,=\,57.5\% (18K checkpoint, single seed s2), suggesting that SDF gradients do not produce a lasting behavioral effect. The per-task CAR ordering is preserved across all conditions and invariant to $\lambda$ (Fig.~\ref{fig:lambda_and_fingerprint}, right), indicating that avoidance patterns reflect the demonstration distribution rather than geometric supervision.

\begin{figure}[t]
\centering
\includegraphics[width=\columnwidth]{fig5_lambda_and_fingerprint}
\caption{(Left) Effect of SDF loss weight $\lambda$ on CAR and TSR ($\lambda = 0$: LoRA-only). All $\lambda$ values remain within the null band ($p > 0.05$). (Right) Per-task CAR fingerprint; the task ordering is preserved across all $\lambda$ values, confirming demonstration-driven avoidance.}
\label{fig:lambda_and_fingerprint}
\end{figure}

\paragraph{CAR--TSR trade-off.}\label{sec:tsr_collapse}
Fine-tuning degrades TSR by ${\sim}$17--22\,pp (71.5\% $\to$ 49.7--54.0\%) due to distributional shift, not learned paralysis (all models retain $>$\,0.37\,m displacement). The Safe Success Rate (SSR)/TSR ratio (SSR\,=\,TSR\,$\cap$\,CAR) increases from 0.39 to 0.51--0.60 (Table~\ref{tab:ssr}): among completing episodes, more avoid collisions after fine-tuning.

\begin{table}[t]
\centering
\caption{SSR and SSR/TSR ratio across conditions. The ratio increase from 0.39 (baseline) to 0.51--0.60 (fine-tuned) indicates that CAR gains reflect learned safer execution, not reduced task exposure.}
\label{tab:ssr}
\vspace{2pt}
\footnotesize
\begin{tabular*}{\columnwidth}{l@{\extracolsep{\fill}}cccc}
\toprule
Condition & TSR & CAR & SSR & SSR/TSR \\
\midrule
Pi0.5 baseline & 71.5 & 28.5 & 28.0 & 0.39 \\
LoRA-only (3 seeds) & 49.7 & 54.83 & 27.7 & 0.56 \\
SDF $\lambda{=}$0.005 & 54.0 & 54.5 & 27.8 & 0.51 \\
SDF $\lambda{=}$0.01 & 53.8 & 56.33 & 30.2 & 0.56 \\
SDF $\lambda{=}$0.1 & 50.2 & 57.00 & 30.2 & 0.60 \\
\bottomrule
\end{tabular*}
\end{table}


\paragraph{Dose-response.}\label{sec:lambda_high}
Increasing $\lambda$ by 20$\times$ to 0.1 yields CAR = 57.00\%\,$\pm$\,2.18 (3~seeds: 58.0, 58.5, 54.5\%), within the null distribution ($p = 0.155$; Fig.~\ref{fig:lambda_and_fingerprint}, left). Alternative null constructions yield consistent non-significance (\S\ref{sec:appendix_null}). With $n = 3$ seeds ($\sigma \approx 2.3$\,pp), 80\% power requires effects $\geq$\,5.6\,pp; the observed ${\sim}$2.2\,pp falls below this minimum detectable effect (MDE). The binding constraint is the structural mismatch between EE monitoring and full-body collision surfaces.

\paragraph{Cross-scene validation.}
To test whether the null SDF result generalizes, we replicate the core comparison on a distinct SafeLIBERO suite (\texttt{spatial\_bowl}, five cluttered obstacles; Table~\ref{tab:spatial_bowl}). SDF ($\lambda = 0.005$) yields CAR of 11.5\% vs.\ 12.0\% for LoRA-only ($\Delta = {-}0.5$\,pp), consistent with absent SDF benefit. The lower CAR (11--12\% vs.\ 54--57\%) reflects greater scene difficulty; both conditions are equally affected.

\begin{table}[t]
\centering
\caption{Cross-scene validation on \texttt{spatial\_bowl} (single training seed, 200 episodes, 30K checkpoint). The $-$0.5\,pp CAR difference is well within noise.}
\label{tab:spatial_bowl}
\vspace{2pt}
\small
\begin{tabular}{lcc}
\toprule
Condition & CAR (\%) & TSR (\%) \\
\midrule
LoRA-only & 12.0 & 49.0 \\
SDF $\lambda{=}0.005$ & 11.5 & 51.5 \\
\bottomrule
\end{tabular}
\end{table}

\paragraph{Evaluation-seed robustness.}
To verify that SDF redundancy is not an artifact of evaluation seed\,=\,7, we re-evaluate LoRA\,s0 and SDF $\lambda{=}0.1$\,s0 across five evaluation seeds each (200 episodes per seed; Table~\ref{tab:eval_seed_main}). LoRA\,s0 achieves mean CAR\,=\,55.4\,$\pm$\,2.2\% across seeds 0, 7, 42, 99, 123; SDF $\lambda{=}0.1$\,s0 achieves 57.9\,$\pm$\,0.9\% across seeds 7, 42, 99, 123, 256. The 2.5\,pp gap is within the LoRA-only seed-to-seed range and non-significant. The broad per-task fingerprint (BBQ consistently lowest; CP and Milk contending for the top) is preserved across all seeds (\S\ref{sec:appendix_eval_seed}), confirming that SDF redundancy holds under diverse evaluation initializations, not just the primary seed.

\begin{table}[t]
\centering
\caption{Eval-seed robustness: CAR (\%) across five evaluation seeds each (200 ep./seed). Ranges overlap; SDF redundancy is consistent across initializations.}
\label{tab:eval_seed_main}
\vspace{2pt}
\footnotesize
\begin{tabular*}{\columnwidth}{l@{\extracolsep{\fill}}ccc}
\toprule
Condition & Seeds & Range & Mean$\pm$Std \\
\midrule
LoRA s0 & 0,7,42,99,123 & 53.0--57.5 & 55.4$\pm$2.2 \\
SDF $\lambda{=}$0.1 s0 & 7,42,99,123,256 & 57.0--59.0 & 57.9$\pm$0.9 \\
\bottomrule
\end{tabular*}
\end{table}

The invariance of per-task avoidance patterns across all conditions and $\lambda$ values (Fig.~\ref{fig:lambda_and_fingerprint}, right) suggests that demonstration distribution is the primary determinant of collision avoidance behavior, motivating a data-centric approach to VLA safety where demonstration curation replaces geometric engineering.

